\section{Arm TrustZone TEE 与漏洞谱系}

\begin{frame}[t,shrink=6]{Arm TrustZone TEE 与漏洞谱系：方向开场}
\DeckKicker{方向开场}
\SlideMeta{证据等级}{方向综述}{来源状态：公开材料综合}
{\large\bfseries\color{Ink} Arm TrustZone TEE 与漏洞谱系 的核心是先界定保护对象、管理边界和证据强度，再用三篇主讲材料串起技术演进。\par}\vspace{1.2mm}
\begin{columns}[T,totalwidth=\textwidth]
  \begin{column}{0.45\textwidth}
    \fcolorbox{Rule}{Paper}{%
  \begin{minipage}[t]{0.97\linewidth}
    \vspace{1.1mm}
    \hspace{1.4mm}\begin{minipage}{0.92\linewidth}
      \PanelTitle{内容说明}
      \scriptsize \begin{itemize}
  \item 把 TrustZone 作为 CCA 前史和边界对比，不把 TrustZone threat model 误写成 CCA。
  \item 固定三篇主讲材料；TrustZone 白皮书、survey 和漏洞 SoK 只支撑历史背景、生态和迁移动机。
  \item 先讲方向痛点和设计轴，再逐篇说明贡献、限制、证据强度和可落地条件。
  \item 对规范、Survey、draft、vendor evidence 保持显式标注，避免把背景材料写成一手系统证明。
\end{itemize}
    \end{minipage}
    \vspace{1mm}
  \end{minipage}}
  \end{column}
  \begin{column}{0.49\textwidth}
    \fcolorbox{Rule}{Panel}{%
  \begin{minipage}[t]{0.97\linewidth}
    \vspace{1.1mm}
    \hspace{1.4mm}\begin{minipage}{0.92\linewidth}
      \PanelTitle{三篇主讲材料的分工}
      \scriptsize {\tiny
\begin{tabularx}{\linewidth}{@{}YYY@{}}
\toprule
\textbf{论文} & \textbf{定位} & \textbf{证据边界} \\
\midrule
Arm TrustZone 白皮书 & 开创/基础入口 & E4 / 基础产业证据 \\
Demystifying Arm TrustZone & 代表性改进 & E2 / Background substrate \\
TrustZone-assisted TEE 漏洞 SoK & 当前 SOTA/补充边界 & E2 / Peer-reviewed SoK \\
\bottomrule
\end{tabularx}
}
    \end{minipage}
    \vspace{1mm}
  \end{minipage}}
  \end{column}
\end{columns}
\SourceFootnote{来源：论文原文、官方规范或公开材料｜证据等级：见本页 badge}
\end{frame}


\begin{frame}[t,shrink=6]{Arm TrustZone 白皮书：内容摘要}
\DeckKicker{内容摘要}
\SlideMeta{证据等级}{E4 / 基础产业证据}{来源状态：本地 PDF 已核验}
{\large\bfseries\color{Ink} 介绍 Secure/Non-secure world、NS bit 和 SoC 级隔离的基本模型\par}\vspace{1.2mm}
\begin{columns}[T,totalwidth=\textwidth]
  \begin{column}{0.45\textwidth}
    \fcolorbox{Rule}{Paper}{%
  \begin{minipage}[t]{0.97\linewidth}
    \vspace{1.1mm}
    \hspace{1.4mm}\begin{minipage}{0.92\linewidth}
      \PanelTitle{内容说明}
      \scriptsize \begin{itemize}
  \item 介绍 Secure/Non-secure world、NS bit 和 SoC 级隔离的基本模型
  \item First-party Arm TrustZone whitepaper is the strongest foundational mechanism source for...
\end{itemize}
    \end{minipage}
    \vspace{1mm}
  \end{minipage}}
  \end{column}
  \begin{column}{0.49\textwidth}
    \fcolorbox{Rule}{Panel}{%
  \begin{minipage}[t]{0.97\linewidth}
    \vspace{1.1mm}
    \hspace{1.4mm}\begin{minipage}{0.92\linewidth}
      \PanelTitle{证据定位}
      \scriptsize {\tiny
\begin{tabularx}{\linewidth}{@{}YY@{}}
\toprule
\textbf{字段} & \textbf{内容} \\
\midrule
论文定位 & 开创/基础入口 \\
材料类型 & 厂商白皮书 \\
证据等级 & E4 / 基础产业证据 \\
来源状态 & 本地 PDF 已核验 \\
\bottomrule
\end{tabularx}
}
    \end{minipage}
    \vspace{1mm}
  \end{minipage}}
  \end{column}
\end{columns}
\SourceFootnote{来源：Arm TrustZone 白皮书｜证据等级：E4 / 基础产业证据}
\end{frame}


\begin{frame}[t,shrink=6]{Arm TrustZone 白皮书：研究背景}
\DeckKicker{研究背景}
\SlideMeta{证据等级}{E4 / 基础产业证据}{来源状态：本地 PDF 已核验}
{\large\bfseries\color{Ink} 移动和嵌入式平台需要把安全服务从普通 OS 中隔离出来\par}\vspace{1.2mm}
\begin{columns}[T,totalwidth=\textwidth]
  \begin{column}{0.45\textwidth}
    \fcolorbox{Rule}{Paper}{%
  \begin{minipage}[t]{0.97\linewidth}
    \vspace{1.1mm}
    \hspace{1.4mm}\begin{minipage}{0.92\linewidth}
      \PanelTitle{内容说明}
      \scriptsize \begin{itemize}
  \item 移动和嵌入式平台需要把安全服务从普通 OS 中隔离出来
\end{itemize}
    \end{minipage}
    \vspace{1mm}
  \end{minipage}}
  \end{column}
  \begin{column}{0.49\textwidth}
    \fcolorbox{Rule}{Panel}{%
  \begin{minipage}[t]{0.97\linewidth}
    \vspace{1.1mm}
    \hspace{1.4mm}\begin{minipage}{0.92\linewidth}
      \PanelTitle{问题边界}
      \scriptsize {\tiny
\begin{tabularx}{\linewidth}{@{}YY@{}}
\toprule
\textbf{维度} & \textbf{说明} \\
\midrule
保护对象 & Arm TrustZone TEE 与漏洞谱系 \\
传统瓶颈 & 把 TrustZone 作为 CCA 前史和边界对比，不把 TrustZone threat model 误写成 CCA。 \\
本文入口 & Arm TrustZone 白皮书 \\
证据限制 & E4 / 基础产业证据 \\
\bottomrule
\end{tabularx}
}
    \end{minipage}
    \vspace{1mm}
  \end{minipage}}
  \end{column}
\end{columns}
\SourceFootnote{来源：Arm TrustZone 白皮书｜证据等级：E4 / 基础产业证据}
\end{frame}


\begin{frame}[t,shrink=6]{Arm TrustZone 白皮书：解决方案}
\DeckKicker{解决方案}
\SlideMeta{证据等级}{E4 / 基础产业证据}{来源状态：本地 PDF 已核验}
{\large\bfseries\color{Ink} 在处理器、互连和外设路径传播安全状态，形成双世界系统\par}\vspace{1.2mm}
\begin{columns}[T,totalwidth=\textwidth]
  \begin{column}{0.45\textwidth}
    \fcolorbox{Rule}{Paper}{%
  \begin{minipage}[t]{0.97\linewidth}
    \vspace{1.1mm}
    \hspace{1.4mm}\begin{minipage}{0.92\linewidth}
      \PanelTitle{内容说明}
      \scriptsize \begin{itemize}
  \item 在处理器、互连和外设路径传播安全状态，形成双世界系统
\end{itemize}
    \end{minipage}
    \vspace{1mm}
  \end{minipage}}
  \end{column}
  \begin{column}{0.49\textwidth}
    \fcolorbox{Rule}{Panel}{%
  \begin{minipage}[t]{0.97\linewidth}
    \vspace{1.1mm}
    \hspace{1.4mm}\begin{minipage}{0.92\linewidth}
      \PanelTitle{核心设计拆解}
      \scriptsize \begin{itemize}
  \item 01. Secure world / Non-secure world
  \item 02. NS bit and SoC security state propagation
  \item 03. Secure monitor and trusted service boundary
\end{itemize}
    \end{minipage}
    \vspace{1mm}
  \end{minipage}}
  \end{column}
\end{columns}
\SourceFootnote{来源：Arm TrustZone 白皮书｜证据等级：E4 / 基础产业证据}
\end{frame}


\begin{frame}[t,shrink=6]{Arm TrustZone 白皮书：实验结果}
\DeckKicker{实验结果}
\SlideMeta{证据等级}{E4 / 基础产业证据}{来源状态：本地 PDF 已核验}
{\large\bfseries\color{Ink} 白皮书/industry evidence，无独立实验\par}\vspace{1.2mm}
\begin{columns}[T,totalwidth=\textwidth]
  \begin{column}{0.45\textwidth}
    \fcolorbox{Rule}{Paper}{%
  \begin{minipage}[t]{0.97\linewidth}
    \vspace{1.1mm}
    \hspace{1.4mm}\begin{minipage}{0.92\linewidth}
      \PanelTitle{内容说明}
      \scriptsize \begin{itemize}
  \item 白皮书/industry evidence，无独立实验
  \item 只支撑 TrustZone 基础模型
\end{itemize}
    \end{minipage}
    \vspace{1mm}
  \end{minipage}}
  \end{column}
  \begin{column}{0.49\textwidth}
    \fcolorbox{Rule}{Panel}{%
  \begin{minipage}[t]{0.97\linewidth}
    \vspace{1.1mm}
    \hspace{1.4mm}\begin{minipage}{0.92\linewidth}
      \PanelTitle{实验/证据边界}
      \scriptsize {\tiny
\begin{tabularx}{\linewidth}{@{}YYY@{}}
\toprule
\textbf{对象} & \textbf{可支撑} & \textbf{边界} \\
\midrule
数据/来源 & Arm TrustZone 白皮书 & 本地 PDF 已核验 \\
Baseline/对照 & 以论文或规范原文为准 & 不跨材料外推 \\
指标/对象 & 只使用当前来源可证明的信息 & E4 / 基础产业证据 \\
使用边界 & 可进入本方向演进图 & 不能替代更强证据 \\
\bottomrule
\end{tabularx}
}
    \end{minipage}
    \vspace{1mm}
  \end{minipage}}
  \end{column}
\end{columns}
\SourceFootnote{来源：Arm TrustZone 白皮书｜证据等级：E4 / 基础产业证据}
\end{frame}


\begin{frame}[t,shrink=6]{Arm TrustZone 白皮书：文章评价}
\DeckKicker{文章评价}
\SlideMeta{证据等级}{E4 / 基础产业证据}{来源状态：本地 PDF 已核验}
{\large\bfseries\color{Ink} 历史基础价值高，能解释 Arm TEE 的双世界起点\par}\vspace{1.2mm}
\begin{columns}[T,totalwidth=\textwidth]
  \begin{column}{0.45\textwidth}
    \fcolorbox{Rule}{Paper}{%
  \begin{minipage}[t]{0.97\linewidth}
    \vspace{1.1mm}
    \hspace{1.4mm}\begin{minipage}{0.92\linewidth}
      \PanelTitle{内容说明}
      \scriptsize \begin{itemize}
  \item 设计优势：历史基础价值高，能解释 Arm TEE 的双世界起点。
  \item 局限性：不代表 CCA/RME/RMM 或 confidential VM threat model。
  \item 商业化潜力：已长期进入移动和嵌入式 SoC 生态；风险在 TCB 大、接口复杂和迁移到 CCA 时边界不同。
  \item 证据边界：E4 / Foundational industry evidence；本地 PDF 已核验。
\end{itemize}
    \end{minipage}
    \vspace{1mm}
  \end{minipage}}
  \end{column}
  \begin{column}{0.49\textwidth}
    \fcolorbox{Rule}{Panel}{%
  \begin{minipage}[t]{0.97\linewidth}
    \vspace{1.1mm}
    \hspace{1.4mm}\begin{minipage}{0.92\linewidth}
      \PanelTitle{文章评价}
      \scriptsize {\tiny
\begin{tabularx}{\linewidth}{@{}YY@{}}
\toprule
\textbf{维度} & \textbf{判断} \\
\midrule
设计优势 & 历史基础价值高，能解释 Arm TEE 的双世界起点。 \\
主要局限 & 不代表 CCA/RME/RMM 或 confidential VM threat model。 \\
商业落地 & 已长期进入移动和嵌入式 SoC 生态；风险在 TCB 大、接口复杂和迁移到 CCA 时边界不同。 \\
讲稿定位 & 开创/基础入口; 证据强度 E4 \\
\bottomrule
\end{tabularx}
}
    \end{minipage}
    \vspace{1mm}
  \end{minipage}}
  \end{column}
\end{columns}
\SourceFootnote{来源：Arm TrustZone 白皮书｜证据等级：E4 / 基础产业证据}
\end{frame}


\begin{frame}[t,shrink=6]{Demystifying Arm TrustZone：内容摘要}
\DeckKicker{内容摘要}
\SlideMeta{证据等级}{E2 / Background substrate}{来源状态：本地 PDF 已核验}
{\large\bfseries\color{Ink} 全面梳理 TrustZone 架构、TEE 软件栈、应用和虚拟化相关工作\par}\vspace{1.2mm}
\begin{columns}[T,totalwidth=\textwidth]
  \begin{column}{0.45\textwidth}
    \fcolorbox{Rule}{Paper}{%
  \begin{minipage}[t]{0.97\linewidth}
    \vspace{1.1mm}
    \hspace{1.4mm}\begin{minipage}{0.92\linewidth}
      \PanelTitle{内容说明}
      \scriptsize \begin{itemize}
  \item 全面梳理 TrustZone 架构、TEE 软件栈、应用和虚拟化相关工作
  \item Systematic survey of TrustZone-assisted TEEs for mature mobile and embedded TEE lineage.
\end{itemize}
    \end{minipage}
    \vspace{1mm}
  \end{minipage}}
  \end{column}
  \begin{column}{0.49\textwidth}
    \fcolorbox{Rule}{Panel}{%
  \begin{minipage}[t]{0.97\linewidth}
    \vspace{1.1mm}
    \hspace{1.4mm}\begin{minipage}{0.92\linewidth}
      \PanelTitle{证据定位}
      \scriptsize {\tiny
\begin{tabularx}{\linewidth}{@{}YY@{}}
\toprule
\textbf{字段} & \textbf{内容} \\
\midrule
论文定位 & 代表性改进 \\
材料类型 & SoK/综述 \\
证据等级 & E2 / Background substrate \\
来源状态 & 本地 PDF 已核验 \\
\bottomrule
\end{tabularx}
}
    \end{minipage}
    \vspace{1mm}
  \end{minipage}}
  \end{column}
\end{columns}
\SourceFootnote{来源：Demystifying Arm TrustZone｜证据等级：E2 / Background substrate}
\end{frame}


\begin{frame}[t,shrink=6]{Demystifying Arm TrustZone：研究背景}
\DeckKicker{研究背景}
\SlideMeta{证据等级}{E2 / Background substrate}{来源状态：本地 PDF 已核验}
{\large\bfseries\color{Ink} TrustZone 长期被广泛部署，但公开系统化资料不足，学术系统和产业 TEE 栈边界不统一\par}\vspace{1.2mm}
\begin{columns}[T,totalwidth=\textwidth]
  \begin{column}{0.45\textwidth}
    \fcolorbox{Rule}{Paper}{%
  \begin{minipage}[t]{0.97\linewidth}
    \vspace{1.1mm}
    \hspace{1.4mm}\begin{minipage}{0.92\linewidth}
      \PanelTitle{内容说明}
      \scriptsize \begin{itemize}
  \item TrustZone 长期被广泛部署，但公开系统化资料不足，学术系统和产业 TEE 栈边界不统一
\end{itemize}
    \end{minipage}
    \vspace{1mm}
  \end{minipage}}
  \end{column}
  \begin{column}{0.49\textwidth}
    \fcolorbox{Rule}{Panel}{%
  \begin{minipage}[t]{0.97\linewidth}
    \vspace{1.1mm}
    \hspace{1.4mm}\begin{minipage}{0.92\linewidth}
      \PanelTitle{问题边界}
      \scriptsize {\tiny
\begin{tabularx}{\linewidth}{@{}YY@{}}
\toprule
\textbf{维度} & \textbf{说明} \\
\midrule
保护对象 & Arm TrustZone TEE 与漏洞谱系 \\
传统瓶颈 & 把 TrustZone 作为 CCA 前史和边界对比，不把 TrustZone threat model 误写成 CCA。 \\
本文入口 & Demystifying Arm TrustZone \\
证据限制 & E2 / Background substrate \\
\bottomrule
\end{tabularx}
}
    \end{minipage}
    \vspace{1mm}
  \end{minipage}}
  \end{column}
\end{columns}
\SourceFootnote{来源：Demystifying Arm TrustZone｜证据等级：E2 / Background substrate}
\end{frame}


\begin{frame}[t,shrink=6]{Demystifying Arm TrustZone：解决方案}
\DeckKicker{解决方案}
\SlideMeta{证据等级}{E2 / Background substrate}{来源状态：本地 PDF 已核验}
{\large\bfseries\color{Ink} 按系统软件、TEE、虚拟化和应用场景分类，给 TrustZone 生态建立 survey map\par}\vspace{1.2mm}
\begin{columns}[T,totalwidth=\textwidth]
  \begin{column}{0.45\textwidth}
    \fcolorbox{Rule}{Paper}{%
  \begin{minipage}[t]{0.97\linewidth}
    \vspace{1.1mm}
    \hspace{1.4mm}\begin{minipage}{0.92\linewidth}
      \PanelTitle{内容说明}
      \scriptsize \begin{itemize}
  \item 按系统软件、TEE、虚拟化和应用场景分类，给 TrustZone 生态建立 survey map
\end{itemize}
    \end{minipage}
    \vspace{1mm}
  \end{minipage}}
  \end{column}
  \begin{column}{0.49\textwidth}
    \fcolorbox{Rule}{Panel}{%
  \begin{minipage}[t]{0.97\linewidth}
    \vspace{1.1mm}
    \hspace{1.4mm}\begin{minipage}{0.92\linewidth}
      \PanelTitle{核心设计拆解}
      \scriptsize \begin{itemize}
  \item 01. TEE OS and trusted applications
  \item 02. Hardware-assisted virtualization use cases
  \item 03. IoT and mobile deployment patterns
\end{itemize}
    \end{minipage}
    \vspace{1mm}
  \end{minipage}}
  \end{column}
\end{columns}
\SourceFootnote{来源：Demystifying Arm TrustZone｜证据等级：E2 / Background substrate}
\end{frame}


\begin{frame}[t,shrink=6]{Demystifying Arm TrustZone：实验结果}
\DeckKicker{实验结果}
\SlideMeta{证据等级}{E2 / Background substrate}{来源状态：本地 PDF 已核验}
{\large\bfseries\color{Ink} 规范/Survey，无新实验；Demystifying Arm TrustZone 只支撑术语、taxonomy、接口语义或证据边界。\par}\vspace{1.2mm}
\begin{columns}[T,totalwidth=\textwidth]
  \begin{column}{0.45\textwidth}
    \fcolorbox{Rule}{Paper}{%
  \begin{minipage}[t]{0.97\linewidth}
    \vspace{1.1mm}
    \hspace{1.4mm}\begin{minipage}{0.92\linewidth}
      \PanelTitle{内容说明}
      \scriptsize \begin{itemize}
  \item 本页不展示独立系统实验，也不把二手归纳写成一手结果。
  \item 可支撑：Systematic survey of TrustZone-assisted TEEs for mature mobile and embedded TEE...
  \item 证据等级：E2 / Background substrate；来源状态：本地 PDF 已核验。
  \item 涉及具体平台、协议状态或数值时，转到原始论文、官方规范或厂商材料核验。
\end{itemize}
    \end{minipage}
    \vspace{1mm}
  \end{minipage}}
  \end{column}
  \begin{column}{0.49\textwidth}
    \fcolorbox{Rule}{Panel}{%
  \begin{minipage}[t]{0.97\linewidth}
    \vspace{1.1mm}
    \hspace{1.4mm}\begin{minipage}{0.92\linewidth}
      \PanelTitle{实验/证据边界}
      \scriptsize {\tiny
\begin{tabularx}{\linewidth}{@{}YYY@{}}
\toprule
\textbf{对象} & \textbf{可支撑} & \textbf{边界} \\
\midrule
数据/来源 & Demystifying Arm TrustZone & 本地 PDF 已核验 \\
Baseline/对照 & 以论文或规范原文为准 & 不跨材料外推 \\
指标/对象 & 只使用当前来源可证明的信息 & E2 / Background substrate \\
使用边界 & 可进入本方向演进图 & 不能替代更强证据 \\
\bottomrule
\end{tabularx}
}
    \end{minipage}
    \vspace{1mm}
  \end{minipage}}
  \end{column}
\end{columns}
\SourceFootnote{来源：Demystifying Arm TrustZone｜证据等级：E2 / Background substrate}
\end{frame}


\begin{frame}[t,shrink=6]{Demystifying Arm TrustZone：文章评价}
\DeckKicker{文章评价}
\SlideMeta{证据等级}{E2 / Background substrate}{来源状态：本地 PDF 已核验}
{\large\bfseries\color{Ink} 评价：Demystifying Arm TrustZone 适合做 taxonomy/规范入口；规范/Survey，无新实验，不能替代一手系统证据。\par}\vspace{1.2mm}
\begin{columns}[T,totalwidth=\textwidth]
  \begin{column}{0.45\textwidth}
    \fcolorbox{Rule}{Paper}{%
  \begin{minipage}[t]{0.97\linewidth}
    \vspace{1.1mm}
    \hspace{1.4mm}\begin{minipage}{0.92\linewidth}
      \PanelTitle{内容说明}
      \scriptsize \begin{itemize}
  \item 设计优势：适合写 TrustZone 背景和生态谱系。
  \item 局限性：对 CCA/RME 不完整，不能直接支撑 Realm 或 confidential VM claim。
  \item 商业化潜力：有助于厂商盘点旧 TEE 生态和迁移成本；风险在 legacy API/TA 兼容性与审计负担。
  \item 规范/Survey，无新实验；具体平台结论转到一手材料核验。
\end{itemize}
    \end{minipage}
    \vspace{1mm}
  \end{minipage}}
  \end{column}
  \begin{column}{0.49\textwidth}
    \fcolorbox{Rule}{Panel}{%
  \begin{minipage}[t]{0.97\linewidth}
    \vspace{1.1mm}
    \hspace{1.4mm}\begin{minipage}{0.92\linewidth}
      \PanelTitle{文章评价}
      \scriptsize {\tiny
\begin{tabularx}{\linewidth}{@{}YY@{}}
\toprule
\textbf{维度} & \textbf{判断} \\
\midrule
设计优势 & 适合写 TrustZone 背景和生态谱系。 \\
主要局限 & 对 CCA/RME 不完整，不能直接支撑 Realm 或 confidential VM claim。 \\
商业落地 & 有助于厂商盘点旧 TEE 生态和迁移成本；风险在 legacy API/TA 兼容性与审计负担。 \\
讲稿定位 & 代表性改进; 证据强度 E2 \\
\bottomrule
\end{tabularx}
}
    \end{minipage}
    \vspace{1mm}
  \end{minipage}}
  \end{column}
\end{columns}
\SourceFootnote{来源：Demystifying Arm TrustZone｜证据等级：E2 / Background substrate}
\end{frame}


\begin{frame}[t,shrink=6]{TrustZone-assisted TEE 漏洞 SoK：内容摘要}
\DeckKicker{内容摘要}
\SlideMeta{证据等级}{E2 / Peer-reviewed SoK}{来源状态：本地 PDF 已核验}
{\large\bfseries\color{Ink} 归纳 TrustZone-assisted TEE 的主流漏洞类型和根因\par}\vspace{1.2mm}
\begin{columns}[T,totalwidth=\textwidth]
  \begin{column}{0.45\textwidth}
    \fcolorbox{Rule}{Paper}{%
  \begin{minipage}[t]{0.97\linewidth}
    \vspace{1.1mm}
    \hspace{1.4mm}\begin{minipage}{0.92\linewidth}
      \PanelTitle{内容说明}
      \scriptsize \begin{itemize}
  \item 归纳 TrustZone-assisted TEE 的主流漏洞类型和根因
  \item S\&P 2020 SoK anchors the TrustZone vulnerability and attack-surface narrative.
\end{itemize}
    \end{minipage}
    \vspace{1mm}
  \end{minipage}}
  \end{column}
  \begin{column}{0.49\textwidth}
    \fcolorbox{Rule}{Panel}{%
  \begin{minipage}[t]{0.97\linewidth}
    \vspace{1.1mm}
    \hspace{1.4mm}\begin{minipage}{0.92\linewidth}
      \PanelTitle{证据定位}
      \scriptsize {\tiny
\begin{tabularx}{\linewidth}{@{}YY@{}}
\toprule
\textbf{字段} & \textbf{内容} \\
\midrule
论文定位 & 当前 SOTA/补充边界 \\
材料类型 & SoK/综述 \\
证据等级 & E2 / Peer-reviewed SoK \\
来源状态 & 本地 PDF 已核验 \\
\bottomrule
\end{tabularx}
}
    \end{minipage}
    \vspace{1mm}
  \end{minipage}}
  \end{column}
\end{columns}
\SourceFootnote{来源：TrustZone-assisted TEE 漏洞 SoK｜证据等级：E2 / Peer-reviewed SoK}
\end{frame}


\begin{frame}[t,shrink=6]{TrustZone-assisted TEE 漏洞 SoK：研究背景}
\DeckKicker{研究背景}
\SlideMeta{证据等级}{E2 / Peer-reviewed SoK}{来源状态：本地 PDF 已核验}
{\large\bfseries\color{Ink} TrustZone TEE 暴露接口复杂、TCB 变大，真实系统反复出现漏洞\par}\vspace{1.2mm}
\begin{columns}[T,totalwidth=\textwidth]
  \begin{column}{0.45\textwidth}
    \fcolorbox{Rule}{Paper}{%
  \begin{minipage}[t]{0.97\linewidth}
    \vspace{1.1mm}
    \hspace{1.4mm}\begin{minipage}{0.92\linewidth}
      \PanelTitle{内容说明}
      \scriptsize \begin{itemize}
  \item TrustZone TEE 暴露接口复杂、TCB 变大，真实系统反复出现漏洞
\end{itemize}
    \end{minipage}
    \vspace{1mm}
  \end{minipage}}
  \end{column}
  \begin{column}{0.49\textwidth}
    \fcolorbox{Rule}{Panel}{%
  \begin{minipage}[t]{0.97\linewidth}
    \vspace{1.1mm}
    \hspace{1.4mm}\begin{minipage}{0.92\linewidth}
      \PanelTitle{问题边界}
      \scriptsize {\tiny
\begin{tabularx}{\linewidth}{@{}YY@{}}
\toprule
\textbf{维度} & \textbf{说明} \\
\midrule
保护对象 & Arm TrustZone TEE 与漏洞谱系 \\
传统瓶颈 & 把 TrustZone 作为 CCA 前史和边界对比，不把 TrustZone threat model 误写成 CCA。 \\
本文入口 & TrustZone-assisted TEE 漏洞 SoK \\
证据限制 & E2 / Peer-reviewed SoK \\
\bottomrule
\end{tabularx}
}
    \end{minipage}
    \vspace{1mm}
  \end{minipage}}
  \end{column}
\end{columns}
\SourceFootnote{来源：TrustZone-assisted TEE 漏洞 SoK｜证据等级：E2 / Peer-reviewed SoK}
\end{frame}


\begin{frame}[t,shrink=6]{TrustZone-assisted TEE 漏洞 SoK：解决方案}
\DeckKicker{解决方案}
\SlideMeta{证据等级}{E2 / Peer-reviewed SoK}{来源状态：本地 PDF 已核验}
{\large\bfseries\color{Ink} 按漏洞面、接口、权限边界和实现缺陷建立 taxonomy，解释为什么后续 CCA 需要更强边界\par}\vspace{1.2mm}
\begin{columns}[T,totalwidth=\textwidth]
  \begin{column}{0.45\textwidth}
    \fcolorbox{Rule}{Paper}{%
  \begin{minipage}[t]{0.97\linewidth}
    \vspace{1.1mm}
    \hspace{1.4mm}\begin{minipage}{0.92\linewidth}
      \PanelTitle{内容说明}
      \scriptsize \begin{itemize}
  \item 按漏洞面、接口、权限边界和实现缺陷建立 taxonomy，解释为什么后续 CCA 需要更强边界
\end{itemize}
    \end{minipage}
    \vspace{1mm}
  \end{minipage}}
  \end{column}
  \begin{column}{0.49\textwidth}
    \fcolorbox{Rule}{Panel}{%
  \begin{minipage}[t]{0.97\linewidth}
    \vspace{1.1mm}
    \hspace{1.4mm}\begin{minipage}{0.92\linewidth}
      \PanelTitle{核心设计拆解}
      \scriptsize \begin{itemize}
  \item 01. Trusted application and TEE OS attack surface
  \item 02. Secure monitor / driver interface risk
  \item 03. Vulnerability taxonomy and root-cause grouping
\end{itemize}
    \end{minipage}
    \vspace{1mm}
  \end{minipage}}
  \end{column}
\end{columns}
\SourceFootnote{来源：TrustZone-assisted TEE 漏洞 SoK｜证据等级：E2 / Peer-reviewed SoK}
\end{frame}


\begin{frame}[t,shrink=6]{TrustZone-assisted TEE 漏洞 SoK：实验结果}
\DeckKicker{实验结果}
\SlideMeta{证据等级}{E2 / Peer-reviewed SoK}{来源状态：本地 PDF 已核验}
{\large\bfseries\color{Ink} 规范/Survey，无新实验；TrustZone-assisted TEE 漏洞 SoK 只支撑术语、taxonomy、接口语义或证据边界。\par}\vspace{1.2mm}
\begin{columns}[T,totalwidth=\textwidth]
  \begin{column}{0.45\textwidth}
    \fcolorbox{Rule}{Paper}{%
  \begin{minipage}[t]{0.97\linewidth}
    \vspace{1.1mm}
    \hspace{1.4mm}\begin{minipage}{0.92\linewidth}
      \PanelTitle{内容说明}
      \scriptsize \begin{itemize}
  \item 本页不展示独立系统实验，也不把二手归纳写成一手结果。
  \item 可支撑：S\&P 2020 SoK anchors the TrustZone vulnerability and attack-surface narrative.
  \item 证据等级：E2 / Peer-reviewed SoK；来源状态：本地 PDF 已核验。
  \item 涉及具体平台、协议状态或数值时，转到原始论文、官方规范或厂商材料核验。
\end{itemize}
    \end{minipage}
    \vspace{1mm}
  \end{minipage}}
  \end{column}
  \begin{column}{0.49\textwidth}
    \fcolorbox{Rule}{Panel}{%
  \begin{minipage}[t]{0.97\linewidth}
    \vspace{1.1mm}
    \hspace{1.4mm}\begin{minipage}{0.92\linewidth}
      \PanelTitle{实验/证据边界}
      \scriptsize {\tiny
\begin{tabularx}{\linewidth}{@{}YYY@{}}
\toprule
\textbf{对象} & \textbf{可支撑} & \textbf{边界} \\
\midrule
数据/来源 & TrustZone-assisted TEE 漏洞 SoK & 本地 PDF 已核验 \\
Baseline/对照 & 以论文或规范原文为准 & 不跨材料外推 \\
指标/对象 & 只使用当前来源可证明的信息 & E2 / Peer-reviewed SoK \\
使用边界 & 可进入本方向演进图 & 不能替代更强证据 \\
\bottomrule
\end{tabularx}
}
    \end{minipage}
    \vspace{1mm}
  \end{minipage}}
  \end{column}
\end{columns}
\SourceFootnote{来源：TrustZone-assisted TEE 漏洞 SoK｜证据等级：E2 / Peer-reviewed SoK}
\end{frame}


\begin{frame}[t,shrink=6]{TrustZone-assisted TEE 漏洞 SoK：文章评价}
\DeckKicker{文章评价}
\SlideMeta{证据等级}{E2 / Peer-reviewed SoK}{来源状态：本地 PDF 已核验}
{\large\bfseries\color{Ink} 评价：TrustZone-assisted TEE 漏洞 SoK 适合做 taxonomy/规范入口；规范/Survey，无新实验，不能替代一手系统证据。\par}\vspace{1.2mm}
\begin{columns}[T,totalwidth=\textwidth]
  \begin{column}{0.45\textwidth}
    \fcolorbox{Rule}{Paper}{%
  \begin{minipage}[t]{0.97\linewidth}
    \vspace{1.1mm}
    \hspace{1.4mm}\begin{minipage}{0.92\linewidth}
      \PanelTitle{内容说明}
      \scriptsize \begin{itemize}
  \item 设计优势：能支撑“为什么需要更强 CCA threat model”的历史动机。
  \item 局限性：不能替代 CCA/RME/RMM 机制分析，也不能直接推出 Realm 安全结论。
  \item 商业化潜力：可用于安全审计和迁移评估；落地条件是缩小 TCB、规范接口并建立持续漏洞治理。
  \item 规范/Survey，无新实验；具体平台结论转到一手材料核验。
\end{itemize}
    \end{minipage}
    \vspace{1mm}
  \end{minipage}}
  \end{column}
  \begin{column}{0.49\textwidth}
    \fcolorbox{Rule}{Panel}{%
  \begin{minipage}[t]{0.97\linewidth}
    \vspace{1.1mm}
    \hspace{1.4mm}\begin{minipage}{0.92\linewidth}
      \PanelTitle{文章评价}
      \scriptsize {\tiny
\begin{tabularx}{\linewidth}{@{}YY@{}}
\toprule
\textbf{维度} & \textbf{判断} \\
\midrule
设计优势 & 能支撑“为什么需要更强 CCA threat model”的历史动机。 \\
主要局限 & 不能替代 CCA/RME/RMM 机制分析，也不能直接推出 Realm 安全结论。 \\
商业落地 & 可用于安全审计和迁移评估；落地条件是缩小 TCB、规范接口并建立持续漏洞治理。 \\
讲稿定位 & 当前 SOTA/补充边界; 证据强度 E2 \\
\bottomrule
\end{tabularx}
}
    \end{minipage}
    \vspace{1mm}
  \end{minipage}}
  \end{column}
\end{columns}
\SourceFootnote{来源：TrustZone-assisted TEE 漏洞 SoK｜证据等级：E2 / Peer-reviewed SoK}
\end{frame}


\begin{frame}[t,shrink=6]{Arm TrustZone TEE 与漏洞谱系：方向总结}
\DeckKicker{方向总结}
\SlideMeta{证据等级}{方向综述}{来源状态：公开材料综合}
{\large\bfseries\color{Ink} Arm TrustZone TEE 与漏洞谱系 的结论是：三篇主讲材料共同给出技术演进，但仍要用证据等级约束每个结论。\par}\vspace{1.2mm}
\begin{columns}[T,totalwidth=\textwidth]
  \begin{column}{0.45\textwidth}
    \fcolorbox{Rule}{Paper}{%
  \begin{minipage}[t]{0.97\linewidth}
    \vspace{1.1mm}
    \hspace{1.4mm}\begin{minipage}{0.92\linewidth}
      \PanelTitle{内容说明}
      \scriptsize \begin{itemize}
  \item Arm TrustZone 白皮书：开创/基础入口，E4 / 基础产业证据。
  \item Demystifying Arm TrustZone：代表性改进，E2 / Background substrate。
  \item TrustZone-assisted TEE 漏洞 SoK：当前 SOTA/补充边界，E2 / Peer-reviewed SoK。
  \item 本方向结论：先用三篇主讲材料建立演进关系，再用证据等级控制机制、实验和商业化结论。
\end{itemize}
    \end{minipage}
    \vspace{1mm}
  \end{minipage}}
  \end{column}
  \begin{column}{0.49\textwidth}
    \fcolorbox{Rule}{Panel}{%
  \begin{minipage}[t]{0.97\linewidth}
    \vspace{1.1mm}
    \hspace{1.4mm}\begin{minipage}{0.92\linewidth}
      \PanelTitle{本方向方法对比}
      \scriptsize {\tiny
\begin{tabularx}{\linewidth}{@{}YYY@{}}
\toprule
\textbf{论文} & \textbf{定位} & \textbf{选择理由} \\
\midrule
Arm TrustZone 白皮书 & 开创/基础入口 & First-party Arm TrustZone whitepaper is the strongest foundational... \\
Demystifying Arm TrustZone & 代表性改进 & Systematic survey of TrustZone-assisted TEEs for mature mobile and... \\
TrustZone-assisted TEE 漏洞 SoK & 当前 SOTA/补充边界 & S\&P 2020 SoK anchors the TrustZone vulnerability and attack-surface... \\
\bottomrule
\end{tabularx}
}
    \end{minipage}
    \vspace{1mm}
  \end{minipage}}
  \end{column}
\end{columns}
\SourceFootnote{来源：论文原文、官方规范或公开材料｜证据等级：见本页 badge}
\end{frame}
